% ============================================================================
% Chapter 2 — Theoretical Foundations
% ============================================================================
\section{Theoretical Foundations}
\label{sec:foundations}

This section assembles the four theoretical building blocks that the formal
model in \S\ref{sec:formulation} integrates: (i) an energetic state variable
with defensible micro-foundations; (ii) career capital as an accumulating,
depreciating stock; (iii) complexity as an endogenous load generated by the
leader's own accumulation; and (iv) the capability-trap lens that motivates a
dynamical---rather than cross-sectional---formalization.

\subsection{Energetic micro-foundations: effort--recovery, not ego depletion}
\label{sec:foundations-energy}

The vitality construct $V(t)$ is grounded in two literatures. First, the
effort--recovery model \citep{meijman1998effort} holds that work effort
produces load reactions that are reversible if recovery opportunities are
adequate, and cumulative if they are not; strain is a property of the
\emph{effort--recovery cycle}, not of effort alone. This is inherently a
statement about dynamics---a flow of depletion balanced against a flow of
restoration---and it maps directly onto the recovery and drain terms of the
vitality equation. Second, subjective vitality
\citep{ryan1997subjective,ryan2000self,nix1999revitalization} provides a
validated, fluctuating, self-reported measure of the energy available to the
self, distinct from arousal and from affect, that responds to autonomous
versus controlled regulation. Recovery research operationalizes the
restoration side \citep{sonnentag2007recovery,sonnentag2022recovery}, and the
organizational literature on human energy documents its behavioral
consequences \citep{quinn2012humanenergy,owens2016relational}.

We deliberately do \emph{not} rest the theory on ego depletion
\citep{baumeister1998ego}. Although early meta-analytic work reported a
medium-sized depletion effect \citep{hagger2010ego}, subsequent bias-corrected
meta-analyses found the evidence consistent with a true effect near zero
\citep{carter2015series}, and a large multisite preregistered replication
failed to detect the paradigmatic effect \citep{vohs2021multisite}; see also
\citet{inzlicht2016running} for a mechanistic reinterpretation. The
effort--recovery foundation requires only that sustained load without adequate
recovery degrades available energy over days-to-months timescales---a claim
with robust support in occupational health research
\citep{demerouti2001job,bakker2017jdr,alarcon2011meta}---and is therefore the
appropriate micro-foundation for a model whose time unit is weeks, not
laboratory minutes.

\subsection{Career capital as a stock: extending human-capital theory}
\label{sec:foundations-capital}

Career capital $C(t)$ aggregates the human, social, and reputational assets a
leader accumulates \citep{becker1993human,day2015leadership}. Two properties
matter for the formalization. First, capital is built through \emph{impactful
deployment}: visible, effective work compounds reputation and network
position, which is why the accumulation term in \S\ref{sec:formulation} is
proportional to realized impact rather than to effort. Second, capital
depreciates: skills obsolesce, networks decay, reputations fade
\citep{march1991exploration}, motivating a standard exponential depreciation
term.

The departure from \citet{becker1993human} is the deployment side. Classical
human-capital theory prices the stock; it does not model the \emph{energy
required to deploy it}. DLVT adds an energetic-deployment constraint: impact
is gated by vitality, so a leader can be capital-rich and impact-poor. We
emphasize that this is an extension, not a refutation---the model reduces to a
standard accumulation--depreciation stock when the energy gate is held open.

\subsection{Complexity as an endogenous load}
\label{sec:foundations-complexity}

The third block is the observation that accumulation itself generates load.
Rising capital brings broader spans of control, more interfaces, more
matters requiring attention, and steeper coordination costs
\citep{garicano2000hierarchies,ocasio1997towards,hambrick2005executive,zhu2022executive}.
The demands literature has long distinguished the level of demands from the
resources available to meet them
\citep{karasek1979job,demerouti2001job,podsakoff2023challenge}; DLVT's
specific move is to make the demand level an \emph{increasing function of the
leader's own accumulated capital}, $O = O_0 + \beta C^{\eta}$, so that
demands are endogenous to success. This is the ``too-much-of-a-good-thing''
logic \citep{pierce2013tmgt} given a dynamic transmission mechanism: the
benefit (capital) and the cost (complexity load) are coupled through a single
state variable rather than juxtaposed cross-sectionally. A scope condition is
declared here and revisited in \S\ref{sec:discussion}: the model treats
complexity as a deterministic image of the leader's own capital, abstracting
from market, peer, and organizational sources of load.

\subsection{Capability traps and the case for a closed-form model}
\label{sec:foundations-traps}

The capability-trap literature
\citep{repenning2001nobody,repenning2002capability,repenning2002simulation,rahmandad2016capability}
shows how organizations that divert effort from capability maintenance to
short-run output can slide into self-reinforcing decline. Two features
characterize those models: they operate at the organizational level, and their
signature dynamic is \emph{bistability}---a healthy basin and a trapped basin
separated by a tipping threshold, with hysteresis under parameter sweeps.

DLVT imports the trap \emph{mechanism} (success-driven load crowding out
regeneration) but arrives at a structurally different conclusion at the
individual level: a \emph{single} globally asymptotically stable low-vitality
attractor (Theorem~\ref{thm:gas}). There is no healthy basin to defend and no
tipping point to avoid; under the baseline kernel family the trap is where
the dynamics go from essentially all initial conditions. The theoretical
payoff of this contrast is precision about interventions: in a bistable world
one large push can move the system across the separatrix, whereas in DLVT
only \emph{parameter} changes---specifically recovery-side changes
(Corollary~\ref{cor:scope})---relocate the attractor itself.

Finally, the choice of a two-dimensional closed-form model over a richer
simulation follows the tradition of \citet{may1976simple},
\citet{murray2002mathematical}, and \citet{strogatz2015nonlinear} in ecology
and applied dynamics, of \citet{pluchino2010peter} in organizational
modeling, and of \citet{cramer2016depression} in formal psychopathology: with
two state variables the entire phase plane can be characterized
analytically---nullclines, equilibria, trapping sets, and Dulac
certificates---so that every qualitative claim is a theorem rather than a
simulation summary. The cost is austerity; the benefit is that the theory can
be wrong in public, in closed form.
